% Generated by scripts/assemble.py; corrections belong in source/corrections.json.
\examhead{112}{安徽大学 2026 年数学分析试卷}{109}
% p110-b17; source PDF page 110
\blockanchor{p110-b17}
\question{1}设  $x_{1}=\frac{1}{5}, x_{n}>0, n=1,2,\cdots$ ，且成立等式  $x_{n+1}^{2}-5x_{n}-6=0$ ，证明：数列  $\{x_{n}\}$  收敛，并求其极限.\par
% p110-b18; source PDF page 110
\blockanchor{p110-b18}
\question{2}计算下列极限:\par
% p110-b19; source PDF page 110
\par\Needspace{6\baselineskip}
\blockanchor{p110-b19}
(1) $\lim_{x\to 0}\frac{\sqrt{1 + 4x} - 2\ln(x + 1) - 1}{e^x - \sin 3x + 2x - 1}.$\par
\nopagebreak[4]
% p110-b20; source PDF page 110
\blockanchor{p110-b20}
\displayformula{(2) \lim _ {n \to \infty} \sum_ {k = 1} ^ {n} \frac {\sin {\frac {\pi}{n}}}{2 + \cos {\frac {k \pi}{n}}}.}
% p110-b21; source PDF page 110
\blockanchor{p110-b21}
\question{3}设广义积分 $\int_{a}^{+\infty} f(x) \, \mathrm{d}x$ 收敛, 且导函数 $f'(x)$ 在 $[a, +\infty)$ 上有界, 证明: $\lim_{x \to +\infty} f(x) = 0$.\par
% p111-b00; source PDF page 111
\blockanchor{p111-b00}
\question{4}设 $f(x)$ 在 $[a, b]$ 上具有连续导数，且 $f(a) = 0$，证明：\par
% p111-b01; source PDF page 111
\blockanchor{p111-b01}
(1) $\int_{a}^{b} f^{2}(x) \, \mathrm{d}x \leq \frac{(b - a)^{2}}{2} \int_{a}^{b} [f'(x)]^{2} \, \mathrm{d}x.$\par
% p111-b02; source PDF page 111
\blockanchor{p111-b02}
(2) $\int_0^{2\pi}f(a\cos x + b\sin x)\mathrm{d}x = \int_0^{2\pi}f(\sqrt{a^2 + b^2}\sin x)\mathrm{d}x.$\par
% p111-b03; source PDF page 111
\blockanchor{p111-b03}
\question{5}设函数 $f(x)$ 在区间 $(-1, 1)$ 上有二阶连续导数，且 $\lim_{x \to 0} \frac{f(x)}{x} = 0, f''(x) \neq 0$，证明：级数 $\sum_{n=2}^{\infty} f\left(\sqrt{\sqrt[n]{2} - 1}\right)$ 发散.\par
% p111-b04; source PDF page 111
\blockanchor{p111-b04}
\question{6}设 $u_{n}(x) (n = 1,2,3,\dots)$. 在 $[a,b]$ 上单调, 且函数项级数 $\sum_{n=1}^{\infty} u_{n}(x)$ 在 $x = a$ 和 $x = b$ 处收敛.\par
% p111-b05; source PDF page 111
\blockanchor{p111-b05}
(1) 证明: $\sum_{n=1}^{\infty} u_n(x)$ 在 $[a, b]$ 上一致收敛.\par
% p111-b06; source PDF page 111
\blockanchor{p111-b06}
(2) 证明: $\sum_{n=2}^{\infty} \ln \left(1 + \frac{x}{n \ln^{3} n}\right)$ 在 $[-a, a]$ 上一致收敛, 其中 $0 < a < 2 \ln^{3} 2$.\par
% p111-b07; source PDF page 111
\par\Needspace{6\baselineskip}
\blockanchor{p111-b07}
\question{7}设一元函数 $y = \varphi(x)$ 在 $x = 0$ 附近有定义, $\varphi(0) = 0$, 且在 $x = 0$ 附近满足 $|\varphi(x)| \leq x^2$, 定义二元函数\par
\nopagebreak[4]
% p111-b08; source PDF page 111
\blockanchor{p111-b08}
\displayformula{f (x, y) = \left\{ \begin{array}{l l} \varphi \left(\frac {x y}{\sqrt {x ^ {2} + y ^ {2}}} \ln \left(1 + \sqrt {x ^ {2} + y ^ {2}}\right)\right), & (x, y) \neq (0, 0); \\ 0, & (x, y) = (0, 0). \end{array} \right.}
% p111-b09; source PDF page 111
\blockanchor{p111-b09}
判断 $f(x,y)$ 在点 $(0,0)$ 处是否可微，并证明.\par
% p111-b10; source PDF page 111
\par\Needspace{6\baselineskip}
\blockanchor{p111-b10}
\question{8}计算曲线积分\par
\nopagebreak[4]
% p111-b11; source PDF page 111
\blockanchor{p111-b11}
\displayformula{\oint_ {L} \left(\frac {- y}{x ^ {2} + y ^ {2}} + 2 y\right) \mathrm{d} x + \left(\frac {x}{x ^ {2} + y ^ {2}} + 3 y\right) \mathrm{d} y.}
% p111-b12; source PDF page 111
\blockanchor{p111-b12}
其中曲线 L 的方程为  $\frac{x^{2}}{9} + \frac{y^{2}}{4} = 1$ ，方向为逆时针方向.\par
% p111-b13; source PDF page 111
\blockanchor{p111-b13}
\question{9}第一卦限中, 在曲面 $\frac{x^2}{a^2} + \frac{y^2}{b^2} + \frac{z^2}{c^2} = 1$ 上求一点使得该点切平面与三个坐标平面所围成的立体体积最小.\par
% p111-b14; source PDF page 111
\par\Needspace{6\baselineskip}
\blockanchor{p111-b14}
\question{10}设区域 $\Omega$ 由分片光滑的封闭曲面 $\Sigma$ 围成, $u(x,y,z)$ 在 $\overline{\Omega}$ 上具有二阶连续偏导数, 且在 $\overline{\Omega}$ 上满足\par
\nopagebreak[4]
% p111-b15; source PDF page 111
\blockanchor{p111-b15}
\displayformula{\frac {\partial^ {2} u}{\partial x ^ {2}} + \frac {\partial^ {2} u}{\partial y ^ {2}} + \frac {\partial^ {2} u}{\partial z ^ {2}} = 0.}
% p111-b16; source PDF page 111
\blockanchor{p111-b16}
(1) 证明: $\iint_{\Sigma} \frac{\partial u}{\partial n} \, \mathrm{d}S = 0$, 其中 $n$ 是曲面 $\Sigma$ 的单位外法向量.\par
% p111-b17; source PDF page 111
\par\Needspace{6\baselineskip}
\blockanchor{p111-b17}
(2) 设 $(x_0, y_0, z_0) \in \Omega$ 为一定点, 证明:\par
\nopagebreak[4]
% p111-b18; source PDF page 111
\blockanchor{p111-b18}
\displayformula{u (x _ {0}, y _ {0}, z _ {0}) = \frac {1}{4 \pi} \iint_ {\Sigma} \left(u \frac {\cos (\widehat {\boldsymbol {r} , \boldsymbol {n}})}{r ^ {2}} + \frac {1}{r} \frac {\partial u}{\partial \boldsymbol {n}}\right) \mathrm{d} S.}
% p111-b19; source PDF page 111
\blockanchor{p111-b19}
其中 $r = (x - x_0, y - y_0, z - z_0)$, $r = \sqrt{(x - x_0)^2 + (y - y_0)^2 + (z - z_0)^2}$.\par
% p111-b20; source PDF page 111
\blockanchor{p111-b20}
\question{11}计算积分 $\int_0^\pi \ln \left(1 + \frac{1}{2}\cos x\right)\mathrm{d}x.$\par
